\chapter{实验协议与统计方法}

\section{实验环境}

正式实验运行于WSL Linux环境，GPU为NVIDIA GeForce RTX 5080 Laptop GPU，显存约16GB。代码使用Python、PyTorch、torchvision、timm、scikit-learn、Pillow、NumPy、pandas、matplotlib和seaborn。依赖由\resultfile{requirements.txt}记录，训练日志保存到\resultfile{logs/}，checkpoint保存到\resultfile{outputs/checkpoints/}。

硬件型号影响速度和部分CUDA算子选择，但不应改变数据协议。所有正式结果保存checkpoint设置、epoch历史、类别顺序和设备信息。公开发布时大权重通过release附件分发，并提供SHA256。

\section{数据职责}

患者级Kermany train用于更新参数，validation用于best checkpoint、配方比较、阈值和温度，locked test用于冻结后的内部评价。RSNA train/validation只在混合和域适配实验中使用；RSNA test用于外部压力测试。NIH test只作补充，不参与选择。

\begin{table}[H]
\centering
\caption{数据集合职责与允许操作}
\begin{tabularx}{\textwidth}{lXXX}
\toprule 集合 & 参数训练 & 超参数/阈值选择 & 结果用途\\\midrule
Kermany grouped train & 是 & 否 & 模型开发\\
Kermany grouped validation & 否 & 是 & checkpoint与校准\\
Kermany grouped locked test & 否 & 否 & 最终内部估计\\
RSNA train & 仅适配实验 & 否 & 跨域训练\\
RSNA validation & 否 & 目标域校准/混合checkpoint & 适配开发\\
RSNA test & 否 & 否 & 外部压力测试\\
NIH test & 否 & 否 & 弱标签敏感性\\\bottomrule
\end{tabularx}
\end{table}

旧项目已多次查看官方Kermany和RSNA test。新患者级Kermany test在本轮重新生成，但本报告编写期间已用于严格基线和候选方案比较，因此它仍是课程设计中的锁定评价，而不是未来论文意义上的完全不可见验证。本文如实记录这一事实。

\section{实验矩阵}

正式矩阵分四层：

\begin{enumerate}
  \item 三模型严格基线：DenseNet121、ConvNeXt-Tiny、ViT-B/16，各seed 42/43/44；
  \item 来源诊断：简单图像统计的Kermany/RSNA来源分类；
  \item 低成本稳健DenseNet：强度扰动+标签平滑，各3 seeds；
  \item 跨域DenseNet：简单混合与域均衡，各3 seeds。
\end{enumerate}

旧官方划分实验、随机初始化、冻结head、最后block解冻和全模型RSNA微调保留为历史探索，不与新版患者级结果混成同一主表。这样避免不同数据协议的数字被误解为公平对照。

\section{评价指标}

混淆矩阵定义TN、FP、FN、TP。主要指标为：
\begin{align}
\mathrm{Accuracy}&=\frac{TP+TN}{TP+TN+FP+FN},\\
\mathrm{Precision}&=\frac{TP}{TP+FP},\\
\mathrm{Recall}&=\frac{TP}{TP+FN},\\
\mathrm{Specificity}&=\frac{TN}{TN+FP},\\
F_1&=\frac{2\cdot\mathrm{Precision}\cdot\mathrm{Recall}}{\mathrm{Precision}+\mathrm{Recall}}.
\end{align}

类别比例不平衡时，accuracy可能被多数类主导，因此把balanced accuracy作为双域方案主要门槛。AUROC描述跨阈值排序；AUPRC补充阳性类precision-recall；Brier和ECE描述概率误差。报告不以“F1最高”替代阴性误报分析。

\section{多随机种子}

每个核心设置运行seed 42、43、44。对非线性深度模型，单次结果可能受到初始化、batch顺序和CUDA算子影响。本文报告三个seed的均值与样本标准差，并为三seed概率平均集成结果计算患者级bootstrap区间。三seed数量仍有限，标准差只能描述本次运行范围，不能等同于总体方差精确估计。

\section{患者级bootstrap}

同一患者多张图像相关，不能把每张图像视为完全独立。bootstrap以subject为抽样单位：从测试患者集合有放回抽取相同数量subject，带入这些subject的全部图像，计算指标；重复2000次，取2.5和97.5百分位作为95\%区间。RSNA每个patientId对应一张处理图，患者级与图像级抽样等价；Kermany则明显不同。

集成概率为三seed对应图像肺炎概率的平均。集成结果用于稳定性总结，不替代单seed可部署模型。若比较两方案，应在同一bootstrap重采样中计算指标差，以利用配对结构。当前主文重点判断区间重叠和效果方向，不把未做多重校正的大量$p$值作为证据。

\section{成功标准}

稳健方案预先采用双域门槛：RSNA balanced accuracy相对严格DenseNet基线至少提高3个百分点，Kermany内部balanced accuracy绝对下降不超过1个百分点，三个seed方向一致。若仅固定阈值指标提高而AUC下降，结论为“操作点部分改善”；只有排序、固定阈值和校准共同改善，才称为全面改善。

复杂域方法的进入条件是低成本方法不足。若简单混合已经明显提高外部AUC和balanced accuracy，优先保留可解释方案，不为增加工作量继续堆方法。这个停止规则避免在同一外部test上反复寻找偶然最高值。

\section{完整性门禁}

正式run前要求：患者交叉为0、跨split精确重复为0、标签冲突为0、manifest哈希与冻结摘要一致。run后要求：评价JSON能由预测CSV复算、checkpoint存在、配置记录完整。最终报告要求摘要、正文、结论和README数字一致。

审计脚本对原Kermany返回FAIL是预期行为，对患者级新划分返回PASS。不能为了让所有命令“绿色”而隐藏原数据问题；机器门禁的职责是阻止错误协议被误当作正式结果。

\section{研究解释规则}

本文遵循以下解释顺序：先看数据协议是否有效，再看效应大小和不确定性，最后讨论模型机制。外部性能下降不能仅凭AUC推断临床失效；内部高分也不能推断临床可用。NIH小样本不作模型排名。Grad-CAM不作定位证明。任何事后分析在正文明确标记。

\section{实验注册表与命名规则}

正式结果文件采用\texttt{方案\_数据集\_split\_模型\_seed}命名。strict表示患者级Kermany-only，robust表示强度扰动加标签平滑，mixed\_simple和mixed\_domain\_balanced表示两种跨域策略。训练摘要、评价JSON、预测CSV和图表使用相同后缀。

一个正式run至少需要：训练摘要含完整settings和history；best checkpoint存在；validation、internal test和RSNA test评价存在；预测CSV行数等于数据集样本数；指标复算一致。缺少任一项时不进入聚合。

错误配置run不会删除，以便追踪问题，但不使用正式前缀。报告在审计部分说明为什么排除。这样既不污染主结果，也不掩盖曾发生的配置问题。

\section{比较机会与测试集治理}

模型开发中最容易忽略的是“看了多少次test”。即使每次都没有把test图像放进反向传播，只要研究者根据test分数决定下一方案，test信息已经影响研究。旧项目多轮观察官方Kermany和RSNA，新版明确把旧矩阵标为探索性。

患者级新test在划分生成后用于三模型、稳健和混合方案。因为用户要求本轮完成全计划，无法再另留一个完全不看的Kermany集合。报告因此不把其称为未来投稿意义的最终盲测。若继续研究，应冻结当前方案并创建新的外部集合，而不是继续切分同一5856张图。

RSNA混合方案的train/val与test按patientId隔离，但研究者已经看到Kermany-only的RSNA test后才设计混合。它是合理的改进实验，却可能对RSNA特定问题适应。第三来源复核仍是必要条件。

\section{bootstrap实现细节}

每次bootstrap从唯一subject列表有放回抽取与原列表等长的subject序列。一个subject被抽中两次时，其全部图像在重采样数据中重复两次；没有抽中的subject不出现。随后用集成概率和固定0.5阈值计算指标。

某次重采样若只含一个类别，AUC无定义，应跳过；当前患者数足够，此情况极少。百分位区间没有作偏差校正，属于直观非参数区间。对于模型差异，最严格做法是在同一次抽样中计算$M_A-M_B$，而不是只比较两个独立区间是否重叠。

报告目前给出各模型区间，并使用预设效果门槛做主判断。未来投稿应补充配对差值区间和多个外部site的分层bootstrap。

\section{计算预算与环境记录}

正式基线共9次训练，稳健方案3次，简单混合3次，域均衡3次，共18次新版训练。每次生成内部和外部预测；另有validation评价、温度拟合、阈值冻结和错误分析。GPU训练串行执行，避免多个run争用显存造成不可控差异。

日志记录epoch级train/validation loss和accuracy。当前没有记录能耗，不能做绿色AI定量比较。不同模型训练时间可以在日志时间戳估计，但未作为性能指标，因为WSL进程与后台负载会影响时间。
